\chapter{Conclusion and Future Work}

\section{Summary of Contributions}

This study has presented a simulation engine for maintaining persona fidelity in LLM-based multi-stakeholder policy deliberation, validated through a 4,080-run benchmark across 20 real-world scenarios. The five principal contributions are:

\begin{enumerate}[leftmargin=*, itemsep=3pt]
    \item \textbf{Persona State Controller with dynamic OCEAN calibration.} A 30-feature behavioral extraction pipeline that estimates OCEAN trait expression at each dialogue turn using relative-to-conversation-baseline normalization, enabling accurate trait estimation even in domain-specific contexts where absolute feature counts are uninformative.

    \item \textbf{4-Slot Candidate Pool with sycophancy detection.} A generation-and-selection architecture that produces four rhetorical variants (integrator, planner, challenger, skeptic) per turn and selects the most persona-consistent candidate using a multi-component scoring function with phase-specific weight modifiers.

    \item \textbf{Commitment Lifecycle for contradiction-free dialogue.} A state management system using soft prompt injection (``Your prior positions:'') that achieves zero contradictions across all 4,080 benchmark runs while maintaining behavioral diversity---demonstrating that declarative state presentation outperforms imperative consistency commands.

    \item \textbf{Structured Action Layer with 12 action types.} An action extraction, validation, and execution pipeline that transforms dialogue into traceable decision-making, achieving 96.8\% action-plan alignment and producing an 87.6\% improvement in outcome fidelity (Cohen's $d = 1.14$) over the naive baseline.

    \item \textbf{Quantitative evidence of the RLHF behavioral ceiling.} Empirical demonstration that RLHF training imposes a structural limit on persona simulation fidelity, manifesting as Neuroticism suppression (\texttt{self\_doubt\_count} $= 0.000$ across all runs), Conscientiousness estimation floor (error $\approx 0.222$), and universal de-escalation bias (100\% of adversarial scenarios simulated as cooperative).
\end{enumerate}


\section{Answers to Research Questions}

The benchmark results provide direct answers to the three research questions motivating this study.

\textbf{RQ1: Can LLMs maintain consistent personality profiles across extended multi-stakeholder deliberation?}

Yes, with engine-level control. The engine achieves persona drift MAE of 0.168 with 95\% CI $[0.167, 0.169]$, representing a 5.7\% improvement over uncontrolled generation (Cohen's $d = -0.94$, large effect). The improvement is concentrated in Openness ($-10.0\%$) and Extraversion ($-16.5\%$), which have strong psycholinguistic signals. However, Neuroticism and Conscientiousness remain resistant to intervention due to RLHF suppression and measurement limitations respectively. The engine demonstrates that personality consistency in multi-agent dialogue is achievable but fundamentally bounded by the behavioral repertoire of the underlying LLM.

\textbf{RQ2: Does structured action tracking improve simulation outcome fidelity?}

Yes, substantially. The engine's outcome fidelity (0.302) exceeds the naive baseline (0.161) by 87.6\%, with Cohen's $d = 1.14$ (large effect) and $p = 0.0004$. The action layer's contribution is evidenced by the high action-plan alignment (0.968) and planned action coverage (0.928), indicating that structured actions create a traceable chain from stated intentions to simulated outcomes. This is the single most impactful component of the engine architecture as measured by effect size.

\textbf{RQ3: How does actor scaling affect persona fidelity?}

Actor scaling is robust. Persona drift increases by only +0.002 MAE from 3-actor (0.171) to 10-actor (0.168) configurations, well within the noise floor ($\pm$0.020). The Persona State Controller's per-actor trait tracking ensures that individual monitoring quality does not degrade with group size. However, role diversity decreases substantially with scale (0.636 to 0.339), indicating that the 12-action vocabulary creates structural overlap at higher actor counts. This suggests that scaling beyond 10 actors would require an expanded action vocabulary.


\section{Future Work}

Five directions emerge from this study's findings.

\textbf{Model-level intervention for RLHF ceiling.} The most impactful improvement would address the RLHF behavioral ceiling directly. Fine-tuning or Direct Preference Optimization (DPO) with training data that includes anxiety expressions, self-doubt, and confrontational disagreement could expand the behavioral repertoire available for Neuroticism and low-Agreeableness simulation. This would require careful balancing to avoid compromising safety properties that RLHF is designed to ensure.

\textbf{RAG knowledge enrichment.} The current engine relies on LLM pre-training knowledge for scenario-specific factual grounding. A Retrieval-Augmented Generation (RAG) pipeline that retrieves and injects entity-specific facts (financial data, regulatory timelines, stakeholder histories) could improve outcome fidelity by providing actors with accurate contextual information rather than relying on potentially outdated or hallucinated training data.

\textbf{Human evaluation study.} The 30-feature extraction pipeline has not been validated against human personality assessments. A controlled study comparing feature-based OCEAN estimates against expert psychologist ratings of the same dialogue would quantify the pipeline's construct validity and identify specific features that over- or under-estimate trait expression.

\textbf{Cross-model comparison.} Evaluating the engine with multiple candidate LLMs (GPT-4, Claude, Gemini, Llama) would clarify which aspects of the RLHF ceiling are universal and which are model-specific. Different RLHF implementations may suppress different traits to different degrees, and the optimal candidate pool might combine models with complementary trait expression capabilities.

\textbf{Interactive deployment interface.} The engine currently operates in batch mode with pre-specified simulation scripts. An interactive interface allowing users to define scenarios in natural language, observe deliberation in real time, and intervene mid-simulation (introducing new information, adding stakeholders) would enable the tool's application as a practical policy analysis instrument.


\section{Closing Remarks}

This study demonstrates that persona fidelity in LLM-based multi-stakeholder deliberation is both achievable and fundamentally bounded. The simulation engine's six components---Persona State Controller, 4-Slot Candidate Pool, Sycophancy Detection, Commitment Lifecycle, Action Layer, and Strategic Disposition---collectively produce statistically significant and practically meaningful improvements over uncontrolled generation across all 20 benchmark scenarios.

However, the discovery of the RLHF behavioral ceiling may be the study's most consequential finding. The ceiling is not a limitation of the engine architecture or the evaluation methodology; it is a property of the underlying LLM, imposed by the optimization process that makes these models useful in the first place. RLHF creates models that are helpful, harmless, and honest---but also models that refuse to express self-doubt, resist generating confrontational disagreement, and default to cooperation even when assigned adversarial roles. For policy deliberation simulation, where authentic conflict and emotional expression are essential, this creates a fundamental tension between model safety and simulation fidelity.

The resolution of this tension---enabling authentic behavioral simulation while preserving safety properties---represents the central challenge for the next generation of persona simulation engines. The present study establishes the empirical baseline against which future progress can be measured.
